\chapter{Red Blood Cell Count (RBC)}

\section*{What Is This Test?}
The red blood cell count (RBC count) measures the number of circulating erythrocytes per unit volume of whole blood, typically expressed as millions of cells per microliter ($\times$ 10\textsuperscript{6}/$\mu$L) \cite{tierney2023current}. Red blood cells (erythrocytes) are biconcave, anucleate cells measuring approximately 7--8 $\mu$m in diameter that function primarily in oxygen and carbon dioxide transport. They are produced in the bone marrow under stimulation by erythropoietin, a hormone secreted predominantly by the renal peritubular interstitial cells in response to tissue hypoxia. Each red blood cell circulates for approximately 120 days before being removed by reticuloendothelial macrophages in the spleen, liver, and bone marrow. The RBC count is a core component of the complete blood count and, together with hemoglobin and hematocrit, provides quantitative assessment of red blood cell mass. It is used in the classification of anemia, evaluation of polycythemia, calculation of red blood cell indices (MCV, MCH, MCHC), and monitoring of hematologic disorders \cite{kaushansky2021williams}. When combined with MCV and MCHC, the RBC count enables differentiation between different types of anemia based on cell size and hemoglobin content.

\section*{Alternative Names}
RBC Count; Erythrocyte Count; Red Blood Cell Number; Red Cell Count; Erythrocyte Concentration

\section*{Principle}
Red blood cells are counted on automated hematology analyzers using the electrical impedance (Coulter) principle or optical light scatter \cite{tietz2012textbook}. In impedance counting, erythrocytes are suspended in a conductive diluent and pass single-file through a small aperture across which a direct current is applied. As each cell passes through, it displaces a volume of conductive diluent equal to its size, generating an electrical resistance pulse proportional to the cell volume. The number of pulses within a defined counting period yields the RBC count. Modern analyzers count 50,000--100,000 red blood cells per sample, giving a coefficient of variation of approximately 1\%. Optical scatter methods (Siemens ADVIA) use laser light diffraction and absorbance measurements to count and characterize each cell. In some instruments, the RBC count, along with MCV (directly measured as mean pulse height), is used to calculate hematocrit (HCT = MCV $\times$ RBC / 10). Manual chamber counting using a Neubauer hemocytometer remains a reference method but is rarely used for clinical purposes, being reserved for situations where automated counts are unreliable (e.g., very low counts, cryoglobulinemia) \cite{bain2017dacie}.

\section*{History}
The first quantitative estimates of red blood cell numbers were made by Karl Vierordt in 1852 using a microscope with a calibrated slide. The hemocytometer, or counting chamber, was developed by Louis-Charles Malassez in 1874 and improved by Karl Thoma and Carl Zeiss in 1878, producing the Thoma-Zeiss counting chamber. The Neubauer improved chamber, introduced in the early 1900s, became the standard for manual RBC counting for much of the 20th century. However, manual counting was labor-intensive and imprecise, with a coefficient of variation of approximately 10--20\%. The revolutionary advancement came in 1953 when Wallace Coulter patented the electrical impedance method for counting cells (Coulter principle). The first commercial Coulter Counter Model A was introduced in 1956 and could count red blood cells, white blood cells, and hemoglobin in minutes with vastly improved precision. This technology transformed hematology from a manual, microscope-based discipline to an automated, quantitative science. Subsequent decades brought multiparameter analyzers capable of generating complete hemograms from a single sample in under one minute.

\section*{Why This Test Is Ordered}
\begin{itemize}
  \item Diagnosis and classification of anemia --- reduced RBC count in conjunction with low hemoglobin and hematocrit confirms anemia; combined with MCV to classify as microcytic, normocytic, or macrocytic
  \item Diagnosis of polycythemia --- elevated RBC count with increased hemoglobin and hematocrit suggests polycythemia vera, secondary polycythemia from chronic hypoxia, or relative polycythemia
  \item Calculation of red blood cell indices --- MCV (mean corpuscular volume), MCH (mean corpuscular hemoglobin), and MCHC (mean corpuscular hemoglobin concentration) all require the RBC count
  \item Monitoring bone marrow function --- serial RBC counts assess erythropoietic activity in chemotherapy patients, after bone marrow transplantation, and during erythropoietin therapy
  \item Evaluation of hemolytic disorders --- decreased RBC count due to premature erythrocyte destruction
  \item Preoperative evaluation --- part of standard preoperative screening
  \item Assessment of nutritional status --- vitamin B12, folate, and iron deficiency all reduce RBC production
  \item Evaluation of symptoms including fatigue, pallor, dyspnea, tachycardia, dizziness, and headache \cite{simel2009rational}
\end{itemize}

\section*{Primary vs Secondary Test}
The red blood cell count is a primary component of the complete blood count and an essential parameter in the evaluation of anemia and polycythemia.

\section*{How This Test Is Performed}
A venous blood sample is drawn into a lavender-top (EDTA) tube from a peripheral vein, usually in the antecubital fossa. The skin is cleansed with antiseptic, a tourniquet is applied briefly, and approximately 2--5 mL of blood is collected. The tube is gently inverted 8--10 times to mix with EDTA anticoagulant, which prevents clotting by chelating calcium. The sample is aspirated into an automated hematology analyzer, which dilutes a precise aliquot of blood, lyses white blood cells (leaving red blood cells intact), and counts the erythrocytes using impedance or optical methods. The entire procedure from venipuncture to result generation takes less than 5 minutes for sample collection and 1--2 minutes for analysis.

\section*{What Preparation Is Needed}
No special preparation is required. Fasting is not necessary for the RBC count. However, patients should avoid strenuous physical activity for 30 minutes before blood collection, as exercise can transiently alter plasma volume and red blood cell concentration. Chronic heavy alcohol consumption should be disclosed, as it can suppress erythropoiesis. Patients should inform their healthcare provider of all medications, particularly chemotherapy agents, erythropoietin, testosterone, and recent blood transfusion.

\section*{Contraindications and Precautions}
\begin{itemize}
  \item No absolute contraindications to blood collection. Factors that may affect RBC count results include: pregnancy (physiological decrease due to plasma volume expansion), living at high altitude (>1,500 m causes compensatory increase), smoking (mild erythrocytosis), dehydration (hemoconcentration produces falsely elevated count), prolonged tourniquet application (>1 minute causes hemoconcentration), recent blood transfusion, hemolyzed specimen (falsely decreases count), and very high white blood cell count (>100,000/$\mu$L may interfere with RBC gate on some analyzers). EDTA is the standard anticoagulant. The sample should be processed within 4 hours at room temperature or 24 hours if refrigerated at 2--8\textdegree{}C. Clotted or hemolyzed specimens must be rejected.
\end{itemize}
\section*{Risks}
Minimal risks associated with venipuncture include: mild pain at the needle insertion site, bruising (ecchymosis) at the puncture site, hematoma formation, lightheadedness or vasovagal syncope, and rarely, infection or nerve injury. These risks are the same as for any routine blood draw.

\section*{What Happens After the Test}
RBC count results are typically available within 30--60 minutes as part of the CBC. The healthcare provider reviews the RBC count together with hemoglobin, hematocrit, and the calculated red cell indices (MCV, MCH, MCHC, RDW) \cite{wallach2022interpretation}. An abnormal RBC count prompts further investigation including reticulocyte count to assess bone marrow response, peripheral blood smear review, and, depending on the clinical picture, iron studies, vitamin B12/folate levels, or hemoglobin electrophoresis. Serial RBC counts may be monitored in patients receiving myelosuppressive chemotherapy, erythropoietin therapy, or blood transfusion.

\section*{Understanding Results}

\subsection*{Below Normal}
\begin{itemize}
  \item \textbf{Iron deficiency anemia:} Reduced RBC count with low MCV, low MCH, elevated RDW, and decreased serum ferritin. The most common cause of anemia globally.
  \item \textbf{Anemia of chronic disease:} Mild to moderate reduction in RBC count, normocytic or mildly microcytic, seen in chronic infection, inflammation, rheumatoid arthritis, malignancy, and chronic kidney disease.
  \item \textbf{Thalassemia:} Reduced RBC count with markedly low MCV, normal or elevated RDW, and abnormal hemoglobin electrophoresis (elevated HbA2 in beta-thal trait, HbH in alpha-thal intermedia).
  \item \textbf{Megaloblastic anemia:} Reduced RBC count with elevated MCV (>100 fL) due to vitamin B12 or folate deficiency causing impaired DNA synthesis and ineffective erythropoiesis.
  \item \textbf{Aplastic anemia:} Markedly reduced RBC count with concurrent leukopenia and thrombocytopenia (pancytopenia) from bone marrow failure.
  \item \textbf{Hemolytic anemia:} Reduced RBC count with elevated reticulocyte count, low haptoglobin, elevated LDH, and unconjugated hyperbilirubinemia.
  \item \textbf{Chronic kidney disease:} Normocytic anemia due to decreased erythropoietin production, correlating with declining GFR.
  \item \textbf{Acute blood loss:} Reduced RBC count after volume redistribution and hemodilution following hemorrhage.
  \item \textbf{Bone marrow infiltration:} Leukemia, lymphoma, myelodysplastic syndrome, or metastatic carcinoma replacing normal marrow.
\end{itemize}

\subsection*{Above Normal}
\begin{itemize}
  \item \textbf{Polycythemia vera:} Primary myeloproliferative neoplasm with JAK2 V617F mutation, autonomous erythroid proliferation, elevated RBC count, hemoglobin, and hematocrit, with low serum erythropoietin.
  \item \textbf{Secondary polycythemia:} Chronic hypoxia from COPD, sleep apnea, cyanotic congenital heart disease, high-altitude residence; all cause elevated erythropoietin with compensatory erythrocytosis.
  \item \textbf{Erythropoietin-secreting tumors:} Renal cell carcinoma, hepatocellular carcinoma, cerebellar hemangioblastoma, pheochromocytoma.
  \item \textbf{Relative polycythemia:} Dehydration, diuretic therapy, severe burns, Gaisb\"{o}ck syndrome (stress polycythemia associated with hypertension and obesity).
  \item \textbf{Exogenous androgens:} Testosterone therapy, anabolic steroids, danazol.
  \item \textbf{Carbon monoxide poisoning:} Chronic smoking or occupational exposure.
  \item \textbf{High-affinity hemoglobin variants:} Hb Chesapeake, Hb Yakima, and other variants.
  \item \textbf{Congenital causes:} Primary familial and congenital polycythemia (EPOR mutations); Chuvash polycythemia (VHL mutations).
\end{itemize}

\section*{Normal Results}
\begin{itemize}
  \item \textbf{Adult males:} 4.35--5.65 $\times$ 10\textsuperscript{6}/$\mu$L (4.35--5.65 $\times$ 10\textsuperscript{12}/L) \cite{fischbach2023manual}
  \item \textbf{Adult females:} 3.92--5.13 $\times$ 10\textsuperscript{6}/$\mu$L (3.92--5.13 $\times$ 10\textsuperscript{12}/L)
  \item \textbf{Children (6--12 years):} 4.0--5.2 $\times$ 10\textsuperscript{6}/$\mu$L
  \item \textbf{Children (1--6 years):} 3.8--5.0 $\times$ 10\textsuperscript{6}/$\mu$L
  \item \textbf{Infants (1--12 months):} 3.5--5.0 $\times$ 10\textsuperscript{6}/$\mu$L
  \item \textbf{Newborns (first week):} 4.0--6.6 $\times$ 10\textsuperscript{6}/$\mu$L (decreases to a physiological nadir at 2--3 months)
  \item \textbf{Pregnant females (third trimester):} 3.0--4.5 $\times$ 10\textsuperscript{6}/$\mu$L (physiological hemodilution)
\end{itemize}

\section*{Follow-up Tests}
\begin{itemize}
  \item \textbf{Complete blood count with differential:} Provides hemoglobin, hematocrit, MCV, MCH, MCHC, RDW, WBC, and platelet count for comprehensive evaluation
  \item \textbf{Peripheral blood smear:} Examination of erythrocyte morphology --- size (anisocytosis), shape (poikilocytosis), color (hypochromia, polychromasia), and inclusions (basophilic stippling, Howell-Jolly bodies, Pappenheimer bodies)
  \item \textbf{Reticulocyte count and reticulocyte production index:} Assesses bone marrow response; distinguishes hypoproliferative from hyperproliferative anemia
  \item \textbf{Iron studies:} Serum iron, ferritin, total iron-binding capacity, transferrin saturation
  \item \textbf{Vitamin B12 and serum folate levels:} When macrocytosis is present
  \item \textbf{Hemoglobin electrophoresis:} Detects hemoglobinopathies (HbS, HbC, thalassemias, HbE)
  \item \textbf{Direct antiglobulin test (Coombs test):} For suspected autoimmune hemolytic anemia
  \item \textbf{Haptoglobin, LDH, indirect bilirubin:} Confirm hemolysis
  \item \textbf{Serum erythropoietin level and JAK2 V617F mutation:} Evaluate polycythemia
  \item \textbf{Bone marrow aspiration and biopsy:} When bone marrow failure, infiltration, or primary hematologic malignancy is suspected
\end{itemize}

\section*{References}
\bibliographystyle{unsrtnat}
\bibliography{references}

\clearpage
\section*{Regulatory Status and Authorization Date}
This chapter describes a test category, not a specific commercial product. FDA, CDSCO, EU IVDR, and MHRA approval or registration dates therefore cannot be assigned without the assay manufacturer, product name, instrument, and intended use. For laboratory-developed tests, an individual agency approval date may not exist. See the Preface for the interpretation of regulatory status.

\clearpage
